\paragraph{PIN cache expiration}
\label{sec:pin-cache-expiration}
The PIN cache holds a PIN in the PD for later verification by a credential.
The PD sends the PIN to the credential in a VERIFY command. Successful
verification allows operations that require the PIN, subject to the
credential's access rules.

Cache Validity limits how long the PD stores the PIN for reuse. When the cache
expires, the PD erases the stored PIN. A credential that has already accepted
the PIN remains PIN verified until the credential's rules clear that state.
Mode \code{0x00} reports both whether the PD has a usable cached PIN and
whether the current credential has accepted PIN verification
(Table~\ref{tab:pivsper-status}).

\textbf{Starting and expiring the cache.}
The PD \textbf{SHALL} start the cache timer when PIN acquisition completes
successfully, using elapsed time unaffected by clock adjustments. Delivery of
the completion reply, status queries, and verification leave the deadline
unchanged. A new successful acquisition starts a new deadline. At expiration,
the PD \textbf{SHALL} erase the cached PIN and mark the cache unavailable.

\textbf{Expiration before VERIFY.}
Immediately before submitting VERIFY to the credential transport,
the PD \textbf{SHALL} check the cache deadline. If the deadline has been
reached, the PD \textbf{SHALL} send no VERIFY and return
\texttt{osdp\_PIVSPER} \code{0x0A} (no cached PIN). The PD \textbf{SHALL}
make the deadline check and submission decision as one indivisible step.
At the deadline, the PIN is expired and the PD \textbf{SHALL NOT} submit VERIFY.

\textbf{Expiration after VERIFY.}
Once VERIFY has been submitted, the PD \textbf{SHALL} finish the exchange even
if the cache expires. The PD reports the credential's verification result,
\code{0x0B} (credential communication failed or timed out), or
\code{0x04} (credential removed), as applicable.
The PD may retain a command copy needed to finish
the exchange under the credential transport rules and \textbf{SHALL} erase
that copy when the exchange ends. Each new VERIFY requires a new ACU request.

\textbf{VERIFY without a confirmed result.}
If communication fails after VERIFY submission, the PD cannot determine
whether the credential accepted the PIN. The PD \textbf{SHALL} clear its
reported PIN-verification state and cached PIN, end any SM session, and require
credential preparation to be restarted before another data or authentication
operation. Restarting preparation restores a known security context after an
exchange whose outcome is unknown. A verified credential reply determines the reported result even if delivery
to the ACU is delayed.

\textbf{Reporting expiration to the ACU.}
While idle, the PD erases an expired cached PIN without sending a reply.
The ACU learns the current
state from a status query or the result of a later verification request.
If Auto Execute is waiting for credential presentation, expiration
\textbf{SHALL} end that request with \code{0x0A} on a later poll.
An accepted verification request receives ACK followed by one terminal
PIVSPER, unless the ACU cancels the request. Expiration after the operation
finishes leaves its pending completion reply unchanged.

Figure~\ref{fig:pin-cache-expiration} shows how expiration affects VERIFY.
Appendix~\ref{ex:pin-expiration} provides packet and timing examples.
\begin{figure}[htbp]
\centering
\begin{tikzpicture}[font=\sffamily,
  panel/.style={draw=OSDPBlue,rounded corners=3pt,fill=OSDPLightGray,
    text width=12cm,align=left,inner sep=9pt},
  flow/.style={->,draw=OSDPBlue,thick}]
\node[panel] (acquire) {\textbf{1. Complete PIN acquisition}\\[3pt]
Start the cache timer. Report acquisition completion on a poll.
Queries and successful verification leave the deadline unchanged.\\[3pt]
\textbf{Expiration while idle:} erase the PIN; send no unsolicited reply.\\
\textbf{Expiration while waiting for presentation:} erase the PIN;
return PIVSPER 0x0A (no cached PIN) on a poll.};
\node[panel,anchor=north] (check) at ([yshift=-0.5cm]acquire.south)
{\textbf{2. Check immediately before submitting VERIFY}\\[3pt]
Has the cache deadline been reached?};
\node[panel,text width=5.3cm,anchor=north] (expired) at ([xshift=-3.3cm,yshift=-1cm]check.south)
{\textbf{Deadline reached}\\[3pt]
Clear the cache. Send no VERIFY. Return PIVSPER 0x0A (no cached PIN) on a poll.};
\node[panel,text width=5.3cm,anchor=north] (finish) at ([xshift=3.3cm,yshift=-1cm]check.south)
{\textbf{3. Finish a submitted exchange}\\[3pt]
Before the deadline, submit VERIFY once. Expiration clears the stored PIN while the exchange finishes.
Return the credential result, or the applicable timeout or removal failure.
Erase the PIN copy used for the command when the exchange ends.};
\node[panel,anchor=north] (state) at ([xshift=-3.3cm,yshift=-0.7cm]finish.south)
{\textbf{4. Report cache and verification status}\\[3pt]
Mode 0x00 reports whether the PD has a cached PIN and whether the
credential has accepted PIN verification.
The cache can be empty while the credential remains PIN verified.};
\draw[flow] (acquire.south) -- (check.north);
\draw[flow] (check.south) -- ++(0,-0.5) -| node[pos=0.75,left,font=\small]{Yes} (expired.north);
\draw[flow] (check.south) -- ++(0,-0.5) -| node[pos=0.75,right,font=\small]{No} (finish.north);
\draw[flow] (finish.south) -- ++(0,-0.35) -| (state.north);
\draw[flow] (expired.south) |- ([yshift=0.35cm]state.north) -- (state.north);
\end{tikzpicture}
\caption{PIN cache lifetime and the VERIFY submission boundary. Idle expiration
is silent. Expiration while waiting for presentation ends the pending request
with PIVSPER 0x0A on a poll.}
\label{fig:pin-cache-expiration}
\end{figure}

